\section{Implementation}
\label{sec:impl}

We implemented \tool in C and C++17.
The complete implementation comprises 4{,}350 lines of code across custom mutators, bandit schedulers, headless shims, and target harnesses.
The software compiles under Clang 18.1.3 on Linux x86\_64 (Ubuntu 24.04).

\paragraph{Module Structure.}
The codebase is organized into three target subdirectories within \texttt{src/}.
First, \texttt{src/fastgltf/} contains the core GLB chunk parser, procedural generator, and bandit scheduler.
Second, \texttt{src/raylib/} provides 37 modular target binaries covering audio, rendering, and models.
Third, \texttt{src/drogon/} includes 19 target configurations covering HTTP parsing, multipart forms, WebSockets, and uploads.

\paragraph{Linker Symbol Wrapping.}
In Clang 18, compiling with \texttt{-fsanitize=fuzzer} provides internal SanitizerCoverage callbacks.
Our coverage tracker links \texttt{\_\_wrap\_\_sanitizer\_cov\_trace\_pc\_guard} to intercept branch events before passing control to \texttt{\_\_real\_\_sanitizer\_cov\_trace\_pc\_guard}.
This decoupling allows the bandit scheduler to read edge events without modifying LLVM libraries.

\paragraph{Headless Shims and Event Loop Binding.}
For \raylib, we configured the software rasterizer (\texttt{rlsw}) to render into a virtual pixel buffer, suppressing X11 and Wayland display server checks.
For \drogon, creating a new \texttt{trantor::EventLoop} per iteration triggered assertion failures when threads retained prior loop state.
We resolved this by capturing the existing thread-local loop via \texttt{trantor::EventLoop::getEventLoopOfCurrentThread()}, enabling synchronous execution cycles without thread recreation.
